\documentclass[11pt]{article}
\usepackage{amsmath, amssymb}
\usepackage{geometry}
\usepackage{array}
\usepackage{booktabs}
\geometry{margin=0.7in}
\setlength{\parindent}{0pt}
\setlength{\parskip}{0.35em}
\renewcommand{\arraystretch}{1.15}

\newcommand{\cheatsheettitle}[2]{%
\begin{center}
{\LARGE #1}\\[0.4em]
{\large #2}
\end{center}
}


\begin{document}
\cheatsheettitle{Algebra II Rational Functions Cheatsheet}{Module 5}

\section*{Definition}
\[
f(x)=\frac{p(x)}{q(x)},\qquad q(x)\ne0
\]

\section*{Domain}
Exclude values that make the original denominator zero.
\[
q(x)=0\Rightarrow x\text{ is not in the domain}
\]

\section*{Holes and Vertical Asymptotes}
Factor first.
\[
\frac{(x-2)(x+5)}{(x-2)(x-3)}
\]
Canceled factor:
\[
x=2\Rightarrow \text{hole}
\]
Remaining denominator zero:
\[
x=3\Rightarrow \text{vertical asymptote}
\]

\section*{Intercepts}
\[
x\text{-intercepts: zeros of the simplified numerator}
\]
\[
y\text{-intercept: }f(0)\text{ if defined}
\]

\section*{Horizontal and Slant Asymptotes}
For
\[
f(x)=\frac{p(x)}{q(x)}
\]
\[
\deg p<\deg q\Rightarrow y=0
\]
\[
\deg p=\deg q\Rightarrow y=\frac{\text{leading coefficient of }p}{\text{leading coefficient of }q}
\]
\[
\deg p>\deg q\Rightarrow \text{use polynomial division}
\]
\[
\deg p=\deg q+1\Rightarrow \text{slant asymptote}
\]

\section*{Rational Equations}
Before clearing denominators, state restrictions.
\[
\frac{A}{B}=\frac{C}{D},\qquad B\ne0,\ D\ne0
\]
Then multiply by the LCD and solve.

\section*{Sign Chart}
Critical values come from zeros and vertical asymptotes. Test one value in each interval to determine positive and negative behavior.

\section*{Quick Checks}
\begin{itemize}
  \item Domain restrictions come from the original denominator.
  \item Holes come from canceled factors.
  \item Vertical asymptotes come from uncanceled denominator factors.
  \item Check solutions to rational equations against restrictions.
\end{itemize}

\end{document}
